\documentclass[11pt,twoside,openright]{memoir}
% Repository typography: EB Garamond, loaded immediately after
% \documentclass per ecosystem invariant XII.  `localtheorems' keeps
% platonic.sty's own theorem environments.
\usepackage[localtheorems]{raeez-math-template}
\usepackage{platonic}
\input{integrated_macros}
\title{\textbf{Calabi--Yau Quantum Groups}\\[0.55em]\Large Factorization, Hall Algebras, Chiral Centres, Quantum Doubles, and Automorphic Characters}
\author{Raeez Lorgat}
\date{26 July 2026}
\hypersetup{pdftitle={Calabi--Yau Quantum Groups},pdfauthor={Raeez Lorgat},pdfsubject={Factorization, Hall Algebras, Chiral Centres, Quantum Doubles, and Automorphic Characters},pdfkeywords={Calabi-Yau category, factorization algebra, cohomological Hall algebra, quantum group, chiral centre, Borcherds product}}
\begin{document}
\frontmatter
\maketitle
\thispagestyle{empty}
\cleardoublepage
\chapter*{Abstract}
\addcontentsline{toc}{chapter}{Abstract}

This monograph gives a single complete language for Calabi--Yau categories, local quantum fields, ordered and associative chiral algebras, Hall and BPS algebras, quantum groups, centres, anomalies, and automorphic products.  Its basic input is an integral Calabi--Yau quantum-group datum containing the categorical, BV, boundary, Hall, Hopf, and character layers together with their comparison morphisms.  The cyclic Deligne and shifted-symplectic theorems provide the universal categorical avatar.  The quantum master equation provides factorization observables.  Boundary restriction provides the ordered algebra and bulk-to-centre map.  Oriented extension correspondences provide Hall multiplication.  A paired triangular Hopf datum provides the completed double and exchange.  The K3 chapters construct the universal action envelope generated by Heisenberg, LLV, and local ADE correspondences.  The $K3\times E$ chapters construct the universal Igusa Hall--Borcherds target and state the exact parity-refined criterion for geometric equivalence.  The final BV and Einstein completion identifies the deformation algebra, field equations, Bianchi identities, propagators, anomalies, and gravitational sectors of the same mixed boundary datum.


\chapter*{The theorem spine}
\addcontentsline{toc}{chapter}{The theorem spine}
The work is organized by the order in which geometry creates algebra:
\[
\begin{aligned}
\text{collision}&\longrightarrow\text{ordered fusion}\longrightarrow\text{centre}\longrightarrow\text{Hall extension}\\
&\longrightarrow\text{quantum double}\longrightarrow\text{character}.
\end{aligned}
\]
Each arrow is constructed in a fixed category.  Each inverse limit uses a proper positive filtration.  Each physical interpretation is the image of a displayed open--closed, Hall, current, trace, or determinant morphism.  The central examples are calculated in generators and relations before their abstract generalization.

\chapter*{Notation and normalization}
\addcontentsline{toc}{chapter}{Notation and normalization}
The ground field has characteristic zero.  Stable categories are presentable and idempotent-complete.  Tensor products preserve colimits separately.  Completions are separated inverse limits of finite weight or proper-height quotients.  The Igusa normalization is
\[
 \Dtheta=64\Dmon,\qquad \ChiTen=\Dmon^2,
\]
with $\Dmon$ monic.  Physical charge is additive; Gram degree is carried through its polarization cocycle.  Bar, factorization, and Hopf coproducts retain distinct notation and distinct geometric meanings.

\tableofcontents
\mainmatter

\part{The invariant architecture}
\chapter{The complete datum}
A Calabi--Yau category, a local BV theory, a Hall correspondence, and a quantum group occupy different mathematical categories.  The theory begins by writing all four objects and every comparison map.

\begin{definition}[Integral Calabi--Yau quantum-group datum]
An integral Calabi--Yau quantum-group datum is
\[
 \mathfrak D=(\Ccal,\theta_{\CY};X,\Lcal,\omega,S,Q;
 b,A_b;\M,\mathcal E,p,q,\varphi,o;
 H^+,H^0,H^-,\Delta,\langle-,-\rangle;\mathcal Z),
\]
where:
\begin{enumerate}
\item $(\Ccal,\theta_{\CY})$ is a smooth or proper Calabi--Yau category with its negative-cyclic class;
\item $(X,\Lcal,\omega,S,Q)$ is an actual cyclic elliptic BV theory and an anomaly-cancelled quantization;
\item $b$ is a boundary generator and $A_b=\RHom(b,b)$ is its ordered factorization endomorphism algebra;
\item $(\M,\mathcal E,p,q,\varphi,o)$ is an oriented Hall moduli problem with coefficient theory $\varphi$;
\item $(H^+,H^0,H^-,\Delta,\langle-,-\rangle)$ is a complete paired triangular Hopf datum;
\item $\mathcal Z$ is a character, determinant, or automorphic section of the constructed algebraic object.
\end{enumerate}
All comparison morphisms are part of the datum and preserve their stated filtration, grading, and operadic structure.
\end{definition}

\chapter{The formal categorical avatar}
A $d$-Calabi--Yau category universally supplies framed $\Etwo$ Hochschild cochains and, on its representable moduli stack, a $(2-d)$-shifted symplectic form.  These two structures determine a formal local avatar prior to the choice of a spacetime field theory.

\begin{definition}[Formal factorization envelope]
Fix a presentable formal-disk operadic category.  For a framed $\Etwo$ algebra $B$, let $\FactEnv_{\mathrm{form}}(B)$ be the free formal factorization object generated by $B$ in that category.  This construction is the relative operadic left adjoint to the formal-disk restriction functor.
\end{definition}

\begin{theorem}[Canonical formal Calabi--Yau avatar]
A Calabi--Yau category $(\Ccal,\theta_{\CY})$ determines
\[
 CC^*(\Ccal),
 \qquad
 \FactEnv_{\mathrm{form}}\bigl(CC^*(\Ccal)\bigr),
 \qquad
 \M_{\Ccal}^{(2-d)\text{-symp}}.
\]
A morphism preserving the Calabi--Yau class induces the corresponding morphisms of framed Hochschild algebras, formal factorization envelopes, and shifted-symplectic moduli whenever the moduli functor is representable.
\end{theorem}
\begin{proof}
The cyclic Deligne theorem gives the framed $\Etwo$ structure.  Relative operadic induction gives the formal envelope.  Pantev--To\"en--Vaqui\'e--Vezzosi transgression gives the shifted-symplectic form.  Naturality follows from functoriality of Hochschild theory, the adjunction, and shifted-symplectic transgression.
\end{proof}

The formal avatar is the universal categorical layer.  A geometric BV theory realizes it on spacetime through a supplied open--closed map.  A Hall moduli problem realizes its extension algebra through a supplied comparison.  The distinction makes every construction composable.

\chapter{The assembly theorem}
\begin{theorem}[Integral assembly]
Every integral Calabi--Yau quantum-group datum determines functorially:
\begin{enumerate}
\item a quantum factorization algebra of observables on $X$;
\item an ordered boundary algebra $A_b$ and a bulk-to-centre morphism
\[
 \Obs_Q\longrightarrow Z_{\ord}(A_b);
\]
\item an associative Hall algebra
\[
 H^+=\bigoplus_\gamma H_*(\M_\gamma,\varphi_\gamma\tensor o_\gamma);
\]
\item a completed double $D(H^+,H^0,H^-)$ with universal $R$-matrix on every convergent continuous module category;
\item a local--wrapped current object after a projection to a curve and an elliptic-degree decomposition;
\item a character $\mathcal Z$ obtained from the chosen trace or automorphic construction.
\end{enumerate}
The construction commutes with morphisms of integral data.
\end{theorem}
\begin{proof}
The quantum master equation gives the observable differential and factorization maps.  Boundary restriction gives $A_b$, and the action on every boundary module gives the centre map.  The two-step extension correspondence and base change give associativity of Hall convolution.  The Hopf pairing gives the completed double and its canonical tensor.  Factorization pushforward and the local--wrapped colour give the curve object.  The selected trace gives the character.  Naturality is inherited at every stage.
\end{proof}

\chapter{The physical experiment}
The mathematics is read from one experiment.  Operators approach along a complex curve and produce a singular chiral coefficient.  Defects approach along an oriented line and compose in order.  Motion in a second topological direction produces exchange.  Closing a defect produces a trace.  Extending one brane by another produces Hall multiplication.  Integrating out a field produces a determinant line.  A Calabi--Yau pairing identifies incoming and outgoing states.  The resulting structures are collision, fusion, braiding, trace, extension, anomaly, and duality.  The assembly theorem records them in their native categories and connects them by explicit maps.

\input{chapters/Volume_I_Ordered_Chiral_Geometry.tex}
\input{chapters/Volume_II_Mixed_HT_Deligne_Theory.tex}
\input{chapters/Volume_IV_Calabi_Yau_Quantum_Groups.tex}
\input{chapters/Volume_III_Igusa_Borcherds_Theory.tex}
\input{chapters/Volume_V_Universal_Chiral_BV_and_Einstein_Completion.tex}

\part{The maximal synthesis}
\chapter{The theorem of the whole construction}
\begin{theorem}[Platonic Calabi--Yau quantum-group theorem]
For every integral Calabi--Yau quantum-group datum $\mathfrak D$, the following diagram is defined and commutes up to its canonical coherent homotopies:
\[
\begin{tikzcd}[column sep=4.1em,row sep=3.1em]
CC^*(\Ccal) \arrow[r,"\FactEnv_{\mathrm{form}}"]
 \arrow[d,"\text{moduli}"'] &
\Obs_Q \arrow[r,"\text{boundary}"] \arrow[d,"\OpenClosed"'] &
A_b \arrow[d,"Z_{\ord}"]\\
\M_{\Ccal}^{(2-d)\text{-symp}}
 \arrow[r,"\Hall"] &
H^+ \arrow[r,"\text{paired double}"] &
D(H^+,H^0,H^-).
\end{tikzcd}
\]
A character functor sends the lower row to the corresponding Hall, PBW, Fock, determinant, or automorphic series.  An open--closed complete boundary identifies the middle vertical map with an equivalence.  A presentation-complete Hall comparison identifies the lower right quantum group with the quantum group of the geometric BPS sector.
\end{theorem}
\begin{proof}
The upper row is the formal-to-geometric realization and boundary restriction.  The left vertical arrow is shifted-symplectic transgression followed by the oriented extension correspondence.  The middle vertical arrow is the open--closed action.  The lower row is Hall convolution followed by the paired Drinfeld double.  Compatibility is the base-change and naturality content of the assembly theorem.  The character statement follows from PBW and trace functoriality.  The two equivalence clauses are the defining recognition criteria for open--closed completeness and presentation completeness.
\end{proof}

\chapter{The universal K3 quantum algebra}
Let $S$ be a K3 surface.  The Nakajima Heisenberg correspondences, the Looijenga--Lunts--Verbitsky action, and the local ADE Yangian correspondences act on distinct geometric towers.  Their common action algebra is defined as the complete algebra generated by their kernels inside the endomorphisms of the direct sum of these towers.

\begin{theorem}[K3 action-envelope theorem]
The completed algebra generated by the Heisenberg, LLV, and local ADE kernels exists as a quotient of their completed amalgamated free product over the common Cartan and grading operators.  It is initial among complete algebras receiving the three geometric actions with their actual mixed commutators.  Its Tannaka coend reconstructs the corresponding coordinate bialgebra whenever the generated module category is rigid and carries a faithful continuous fibre functor.
\end{theorem}
\begin{proof}
Complete filtered algebras admit degreewise amalgamated free products.  Quotient by the kernel of the combined action map to the endomorphism algebra of the geometric tower.  This quotient has the asserted universal property.  Under rigidity and a faithful fibre functor, the coend construction gives the coordinate bialgebra and recovers the action category.
\end{proof}

This object gives the phrase ``K3 quantum algebra'' a precise universal meaning.  It preserves the three established geometric sources and stores their mixed commutators instead of identifying them by rank.

\chapter{The central compact example}
For $X=K3\times E$, the universal Igusa object supplies the algebraic Hall--Borcherds target, its quantum group, its local--wrapped current object, and its automorphic line.  A reduced oriented compact Hall source, when supplied, is compared with this target by the geometric recognition theorem.  Thus the formal, algebraic, analytic, and automorphic layers are closed; the compact geometric content is concentrated in a single explicit realization datum.

\begin{theorem}[Universal--geometric equivalence criterion]
Let $\GeoDatum$ be a reduced oriented $K3\times E$ Hall datum with polarized Gram grading.  A comparison
\[
 \Phi_{\mathrm{BPS}}:\Prim(\GeoDatum)\longrightarrow\gIg_+
\]
extends to an equivalence of completed doubles, current factorization objects, centres, and determinant lines exactly when:
\begin{enumerate}
\item it maps the geometric simple states to the real and imaginary simple generators;
\item it is surjective on primitives;
\item its parity-refined PBW character equals the Igusa root character;
\item it intertwines the coproduct, Hopf pairing, local--wrapped actions, and orientation line.
\end{enumerate}
\end{theorem}
\begin{proof}
The first condition gives a presentation morphism.  The second and third give a primitive isomorphism by minimal-height PBW recognition.  The fourth transports the additional structures.  Completion is exact in every proper-height degree, hence the equivalence passes to the inverse limit.
\end{proof}

\chapter{Coda: geometry becomes algebra and preserves its origin}
A Calabi--Yau category supplies duality.  A BV theory supplies locality.  A boundary supplies order.  A second transverse direction supplies exchange.  A moduli stack supplies states.  Extension correspondences supply multiplication.  A pairing supplies the double.  A trace supplies a character.  A Borcherds lift supplies an automorphic denominator.  The integrated datum preserves the source of every operation.  This preservation is the simplicity of the subject: each geometric motion has one algebraic image, and the whole theory is their functorial composition.

\backmatter
\bibliographystyle{alpha}
\bibliography{references}
\end{document}
